\documentclass[11pt,a4paper]{article}
\usepackage[T1]{fontenc}
\usepackage[utf8]{inputenc}
\usepackage[czech]{babel}
\usepackage{lmodern,microtype}
\usepackage{amsmath,amssymb,mathtools}
\usepackage{geometry}
\geometry{left=19mm,right=19mm,top=21mm,bottom=21mm}
\usepackage[dvipsnames]{xcolor}
\usepackage{enumitem,tabularx,booktabs}
\usepackage[most]{tcolorbox}
\usepackage{fancyhdr,lastpage}
\usepackage[hidelinks]{hyperref}
\definecolor{deep}{HTML}{203955}
\definecolor{teal}{HTML}{176E76}
\definecolor{slate}{HTML}{506173}
\definecolor{pale}{HTML}{F3F7F9}
\setlength{\parindent}{0pt}
\setlength{\parskip}{5pt}
\setlist[enumerate]{leftmargin=6mm,itemsep=5pt,topsep=3pt}
\pagestyle{fancy}\fancyhf{}
\fancyhead[L]{\sffamily\footnotesize\color{slate} MAT · 4. ročník · kombinatorika}
\fancyhead[R]{\sffamily\footnotesize\color{slate} Hodina \LESSONNUM\, / 112}
\fancyfoot[L]{\sffamily\footnotesize\color{slate} Výukový podklad · samostatný dokument}
\fancyfoot[R]{\sffamily\footnotesize\color{slate} \thepage/\pageref{LastPage}}
\renewcommand{\headrulewidth}{0.35pt}
\renewcommand{\footrulewidth}{0pt}
\newcommand{\sect}[1]{\par\vspace{2mm}{\Large\sffamily\bfseries\color{deep} #1}\par\vspace{-0.5mm}\noindent\color{teal}\rule{\linewidth}{0.7pt}\color{black}\vspace{1mm}}
\tcbset{theobox/.style={enhanced,breakable,colback=pale,colframe=teal!45!white,boxrule=.6pt,arc=1mm,left=3mm,right=3mm,top=2mm,bottom=2mm}}
\newcommand{\LESSONNUM}{39}
\begin{document}

{\Huge\sffamily\bfseries\color{deep} Binomická věta a koeficienty členů}\par
{\sffamily\large\color{teal} Výukový podklad · 39. vyučovací hodina}\par
\textbf{Hodinová dotace:} 45 minut\quad \textbf{Ročník:} 4. SPŠ\quad \textbf{Rozpis:} TP 2.3\par
\begin{tcolorbox}[theobox]\textbf{Výukový cíl.} Dokázat význam binomických koeficientů a rozvíjet \((a+b)^n\) včetně záporného členu a vyhledávání koeficientů.\end{tcolorbox}
\sect{Teorie a vysvětlení}
Pro každé \(n\in\mathbb N_0\) platí \textbf{binomická věta}
\[(a+b)^n=\sum_{k=0}^n\binom nk a^{n-k}b^k.\]
V každém z \(n\) činitelů \((a+b)\) volíme \(a\), nebo \(b\). Aby vznikl člen \(a^{n-k}b^k\), musíme vybrat \(b\) přesně v \(k\) závorkách, což lze \(\binom nk\) způsoby.\par
\textbf{Pravidlo práce:} mocnina prvního členu klesá od \(n\) k nule, mocnina druhého roste od nuly k \(n\), koeficienty jsou v řádku \(n\) Pascalova trojúhelníku. Pro rozdíl dosadíme \(b=-c\).
\sect{Řešené příklady}
\textbf{Příklad 1.} Rozviň \((x+2)^4\).\par
\textit{Řešení.} Koeficienty \(1,4,6,4,1\). Členy \(x^4\), \(4x^3\cdot2\), \(6x^2\cdot2^2\), \(4x\cdot2^3\), \(2^4\). Dostáváme \(x^4+8x^3+24x^2+32x+16\).\par\smallskip
\textbf{Příklad 2.} Rozviň \((2x-y)^3\).\par
\textit{Řešení.} Položíme \(a=2x, b=-y\) a použijeme koeficienty \(1,3,3,1\): \(8x^3-12x^2y+6xy^2-y^3\).\par\smallskip
\textbf{Příklad 3.} Najdi koeficient u \(x^3\) v \((x+2)^5\).\par
\textit{Řešení.} Obecný člen je \(\binom5k x^{5-k}2^k\). Pro \(x^3\) potřebujeme \(5-k=3\), tedy \(k=2\). Koeficient \(\binom52\cdot2^2=10\cdot4=40\).\par\smallskip
\textbf{Příklad 4.} Najdi koeficient u \(x^4\) v \((2x-1)^6\).\par
\textit{Řešení.} Obecný člen \(\binom6k (2x)^{6-k}(-1)^k\). Pro \(x^4\): \(k=2\). Koeficient \(\binom62\cdot2^4\cdot(-1)^2=15\cdot16=240\).\par\smallskip
\textbf{Příklad 5.} Použij binomickou větu ke zdůvodnění \(\sum_{k=0}^n\binom nk=2^n\).\par
\textit{Řešení.} Vezmeme \(a=b=1\): \((1+1)^n=\sum_{k=0}^n\binom nk\cdot1^{n-k}\cdot1^k=\sum_{k=0}^n\binom nk\).\par\smallskip
\clearpage\sect{Procvičování pro studenty}
\begin{enumerate}
\item Rozviň \((a+b)^3\).
\item Rozviň \((x-1)^4\).
\item Najdi koeficient u \(x^2\) v \((x+3)^4\).
\item Najdi konstantní člen v \((2x+5)^3\).
\item Najdi koeficient u \(x^3\) v \((1-2x)^5\).
\end{enumerate}
\sect{Řešení a zdůvodnění k procvičování}
\begin{enumerate}
\item \(a^3+3a^2b+3ab^2+b^3\).
\item \(x^4-4x^3+6x^2-4x+1\).
\item \(\binom42 3^2=6\cdot9=54\).
\item Vzniká volbou \(5\) ve všech třech závorkách: \(5^3=125\).
\item \(\binom53(-2)^3=10\cdot(-8)=-80\).
\end{enumerate}
\sect{Metodické poznámky}
\textbf{Typické chyby:}
\begin{itemize}[leftmargin=5mm,itemsep=1pt]
\item Mocnina \(x\) určuje vhodný index \(k\); nelze číst přímo pouze koeficient z Pascalova trojúhelníku.
\item U \((a-b)^n\) se znaménka střídají; záporné \(b\) umocňujeme v každém členu.
\item Rozvoj má přesně \(n+1\) členů před případným slučováním.
\end{itemize}
\textbf{Doporučený průběh 45 minut:} 0–12 odvození; 12–25 rozvoje; 25–34 hledání koeficientu; 34–42 procvičení; 42–45 shrnutí.\par
\textbf{Návaznost:} Příští hodina obsahuje krátké ověření variací, permutací a kombinací; binomickou větu upevníme v souhrnných úlohách.\par
\vspace{2mm}\noindent{\footnotesize\color{slate} \textbf{Východiska:} rozpis výuky MAT, 4. ročník, hodiny 29--41; tematický plán SPŠ; srovnání tematického rozsahu a náročnosti s realisticky.cz, oddíl Kombinatorika (lekce 090101--090119). Výklad, příklady a řešení jsou nově autorsky sestavené, nikoli převzatá sbírka.}\par
\end{document}